\def\week{85}
\def\topic{Drawing machines}
\pdfinfoomitdate=1
\pdftrailerid{}
\pdfsuppressptexinfo=15
\documentclass[11pt,letterpaper]{article}
\usepackage[margin=.68in,headheight=16pt,headsep=18pt,footskip=24pt]{geometry}
\usepackage[T1]{fontenc}
\usepackage{lmodern,amsmath,amssymb,tikz,fancyhdr,array}
\usepackage[hidelinks]{hyperref}
\pagestyle{fancy}\fancyhf{}
\fancyhead[L]{\small Week \week / \topic / Facilitator}
\fancyfoot[L]{\footnotesize Bellingham Math Circle / Week \week / W\week-F-v1}
\fancyfoot[R]{\footnotesize\thepage}
\renewcommand{\headrulewidth}{.4pt}
\setlength{\parindent}{0pt}\setlength{\parskip}{7pt}
\newcommand{\heading}[1]{\par\vspace{4pt}{\large\bfseries #1}\par}
\newcommand{\sub}[1]{\par\textbf{#1}\par}
\begin{document}
\heading{Mathematical overview}
In the ideal planar machine, $O,X,Q$ lie on one rigid bar with $OX=XQ=L$; $Q,Y,P$ lie on another with $QY=YP=L$. Short bars $XT$ and $TY$ have length $L$. On the simple, noncrossed, noncollapsed parallelogram branch $XQYT$, the tracer $T$ and output $P$ satisfy
\[
 P=O+2(T-O).
\]
This is a dilation of factor 2 about the fixed anchor $O$. It holds throughout the permitted motion, not just at a few measured poses. With $O$ as origin, take directed half-bars $u$ and $v$. Then $X=u$, $Q=2u$, $Y=2u+v$, $T=u+v$, and $P=2u+2v$. Adding directed strips gives a concrete explanation; formal vector notation is optional for children.

For two tracer positions, $P_2-P_1=2(T_2-T_1)$. Thus all corresponding lengths double, angles and orientation stay the same, and area is multiplied by four. A square of twice the side contains four equal original squares; grid-square and half-square dissections establish the same fact for the packet's polygons. This is not a claim of exact measured copies from an imperfect homemade linkage.

Swapping controlled and drawing points gives the inverse scale $1/2$ and area factor $1/4$. If the tracer picture stays fixed while the anchor moves from $O$ to $O+d$, the new copy is $P'=2T-(O+d)=P-d$: it moves the same distance in the opposite direction. This comparison requires the same physical input picture and a viable pose for both anchors.

\sub{What children discover}
Concrete copying and inverse control are an approximate grades 2--3 entry. Upper readiness supports the linkage explanation, area comparison and moving-center investigation. There is no independent K--1 edition: controlled tracing and mechanism recovery are real prerequisites. Adult reading or scribing is welcome; an adult steering every movement is not independent investigation.

Copies and measurements are experiments. A pattern suggested by several poses is a conjecture. The rigid-bar/parallelogram argument explains every permitted pose. Inaccurate copies may expose loose joints or off-axis pen tips, rather than a counterexample to the ideal theorem.

\sub{Limits and status}
Prepared, independently reviewed digital prototype; unpiloted. Physical assembly, coaxial mounts, clearance, hand access and print fit have not been rehearsed. Do not use a crossed or collapsed linkage, and do not describe this packet as hardware-ready until a complete working kit passes the rehearsal on the next page. No statement assumes arbitrary four-bar assemblies have the dilation property.
\newpage
\heading{Preparation and the first visit}
Use three pair machines if both older tables participate: two for the four third graders, one for the three older children rotating tracer, observer and referee. The anchored K--1 adult should use a familiar independent activity; a short supported station is optional only after the hardware succeeds. Each pair needs an assembled machine, support board, joined mat, removable tracing overlays, ruler, straightedge and pencils. Adults manufacture and reset the machine.

\sub{Hardware that must be real}
Two rigid long bars have hole centers at 0,80,160 mm; two short bars have end-hole centers 80 mm apart. Attach $O$ to a fixed rotating anchor; pivot at $X,Q,Y$; join the two short bars at a freely turning coaxial tracer hinge $T$; put the output pencil tip directly on the $P$ axis. A pencil taped beside a bar invalidates the exact scale model. Smooth washers/spacers prevent binding. Choose actual sleeve diameters to fit the actual pencil and hardware; printed circle centers are not a complete manufacturing specification.

Use a flat nonslip support at least $450\times400$ mm. Join the two Letter mat sheets at actual size, matching their registrations and checking the 100 mm segment. Center the joined mat on the support: linkage joints extend outside the paper. Preserve full-zone support and hand clearance. Put fresh overlays on the same registration, retaining the fixed anchor and earlier copies.

\sub{Required physical rehearsal (not yet performed)}
Assemble one machine. Confirm a simple open rhombus and the midpoint topology. Put $T$ at nine positions spanning the marked zone and check $O,T,P$ alignment and $OP=2OT$. Trace straight, cornered and closed paths; exchange drawing/control roles; test the alternative anchor if that route will be used. Check that neither point lifts, drags or crosses branches, that all bars stay supported, and that two children's hands fit. Reject loose or easily crossing assemblies. Mark the intended branch; an adult resets any stuck/collapsed assembly. A screen animation does not satisfy this rehearsal.

Budget 60--120 minutes for a first prototype, potentially longer if mounts must be fabricated; then roughly 15--25 minutes per additional proven kit and 10--15 minutes for printing/joining mats. These are untested budgeting estimates. If the mechanism fails, postpone its student tasks and keep the mathematics for a later visit.

\sub{Launch together}
After the customary five-minute run, let children handle the mechanism gently with tips raised. Anchor $O$ and trace one small curved stroke with $T$, pausing midway to show the emerging $P$ stroke. Keep the same anchor and input through the completed demonstration. Say: ``You choose the tracer's path. The other pencil makes its copy. Stop if the bars get stuck.'' Give each child a turn before distributing tasks. Leave scale and area discoveries for their attempts.
\newpage
\heading{Shared packet: copying and control}
The shared packet has four problems. A first visit can explore Problems 1--2 for most of the table time, then compare copies. Problems 3--4 can follow when children want a rule or another control challenge. Give only the inserts for the current task and keep one useful overlay for the next visit. A separate upper packet has its own Problems 1--4.

\sub{Shared Problem 1: paths A, B, C}
Every stroke of the output is the same shape and orientation at twice the scale, positioned by dilation about $O$. In A, 35 mm then 20 mm becomes 70 mm then 40 mm. In B, the triangle's corresponding sides double; it does not rotate or turn over. In C, the staircase segment lengths 10,15,20,25,10 mm become 20,30,40,50,20 mm. Corners and the order of the path are retained. Several repeated trials are useful, especially if corners introduce pen drag.

Hold back measuring instructions until copies exist. If useful, offer one strip to compare corresponding segments or ask which tip drew the larger picture. An apparent rotation, reflection or varying scale prompts an anchor/branch/mount check, not an immediate correction of the child's observation. The ideal claim and the real machine's precision are different questions.

\sub{Shared Problem 2: visit D and trace E}
Controlling T while watching P is an inverse task. With $O=(40,140)$ in mat millimeters, the inverse is $T=(P+O)/2$. The five D target dots correspond to tracer positions $(115,125)$, $(130,160)$, $(155,150)$, $(150,120)$, $(125,140)$. Children may visit them in any order; the task is control, not a prescribed route.

E's house vertices correspond to $(120,125)$, $(150,125)$, $(150,145)$, $(135,157.5)$, $(120,145)$, then return. Children need not use coordinates or fractions: the two tips trace both paths physically. Give a fresh overlay when the first path becomes confusing. A held hint is to place T tentatively and observe which direction P moves, preserving the fixed anchor.

\sub{Shared Problem 3: a rule for poses}
No permitted ideal pose breaks $O,T,P$ collinearity or $OT=TP$ (equivalently $OP=2OT$). Children can make experiments with a straightedge, string or reusable strip if they ask how to compare. Finite testing suggests a conjecture; the upper packet's bar argument establishes it for every permitted pose. Do not count a stuck, crossed or off-axis-tip assembly as a counterexample to these assumptions.

\sub{Shared Problem 4: stay inside F}
F is a $40\times30$ mm output rectangle. Its inverse is a $20\times15$ mm tracer rectangle, with vertices $(130,132.5)$, $(150,132.5)$, $(150,147.5)$, $(130,147.5)$. All paths within that inverse region produce paths inside F. A tracer rectangle whose copy fills F's outline is therefore exactly this inverse rectangle; there is no larger admissible rectangle by inclusion or area. This answer is a geometric control constraint, not a largest-perimeter contest. Children can discover its boundary by controlling P onto F's sides and retaining T's marks.
\newpage
\heading{Upper packet: explanation and return visits}
\sub{Upper Problem 1: why doubling persists}
Name the equal directed half-bars $u=\overrightarrow{OX}=\overrightarrow{XQ}$ and $v=\overrightarrow{QY}=\overrightarrow{YP}$. In the intended parallelogram, $\overrightarrow{XT}=v$ and $\overrightarrow{TY}=u$. Thus $\overrightarrow{OT}=u+v$ and $\overrightarrow{OP}=2u+2v=2(u+v)$. Loose directed paper strips or arrows on the diagram can show this without symbolic algebra. A page of measurements cannot replace that rigid-frame explanation. Equal-length marks alone do not specify the operating branch: keep the same simple rhombus as the real assembly.

\sub{Upper Problem 2: area of G, H, I}
The 10 mm square is one area unit. G has area 4 units and its copy 16; H also has area 4 and its copy 16; I is the right-isosceles triangle with 30 mm legs, area $4\frac12$, and its copy 18. These are square units, not boundary lengths or raw piece counts. Two half-square triangles make one unit. The triangle's diagonal aligns with the supplied half-square pieces, so exact covering is possible in the ideal drawing. Cover one shape/copy at a time and reuse the kit; it is not necessary to leave all six covered simultaneously.

A held hint is to place four unit squares into a doubled unit square. A grid-square dissection of G and H then gives four copies of every original unit. For I, pair two congruent triangles into a square, use the fourfold square area, then take half. This explains all three and extends to polygons decomposed into these units. Physical tracing thickness and play in joints may prevent exact edge matching; compare the intended geometry with the measured result.

\sub{Upper Problem 3: reverse the machine}
An adult swaps control/drawing roles: children move P and T makes the reduced drawing. J's $60\times40$ mm rectangle becomes $30\times20$ mm. Every corresponding arrow length in K halves, retaining angles and orientation. The full relation is now $T=O+\frac12(P-O)$. To make T fill L, whose bounds are $(125,130)$ to $(155,150)$, P must trace the $60\times40$ mm rectangle from $(210,120)$ to $(270,160)$. If the mount does not permit a reliable exchange of roles, postpone this visit rather than silently moving an off-axis pencil.

\sub{Upper Problem 4: change the anchor}
The printed second anchor $O'=(60,140)$ is 20 mm to the right of $O$. Keeping M fixed moves its copy 20 mm left: $P'=2T-O'=P-(20,0)$. To make the second copy land on the first, move every point of M 10 mm right, because $2(T+(10,0))-O'=2T-O$. Its new tracer vertices lie within the intended zone. The shape and area do not change. The demonstration must retain M's physical location while the anchor changes for the first comparison; shifting M at the same time would answer a different question.

The general compensation is half the anchor displacement in the same direction. Ask for a new displacement only after the first two copies are visible; no simultaneous coordinate table is needed. Retain both overlays and the anchor labels as a recoverable record.
\newpage
\heading{Print selection, provenance and use record}
The shared and upper packets are independent consecutively numbered routes. Use the matching lettered materials, not rescaled thumbnails, on the joined mat. The materials include adult templates and setup, the two mat halves, seam bridge, inserts A--M, and area pieces. The source README gives the exact current material page map. Print the drilling templates only for adult fabrication; children use assembled inspected kits.

\sub{Common failure checks}
If an insert seems to demand a stuck pose, stop. Confirm the named insert's registration, correct control point, anchor position and intended branch before trying again. Do not force a joint. Tip offsets or loose pivots can enlarge errors even when bar-hole distances are right. Joined paper may still be too small to support the outer joints; the larger board is required. Adults should not operate both tips to produce an attractive result on a child's behalf.

\sub{Sources and adaptation}
Glen Whitney, \emph{Math Monday: Linkages -- The World ``B.X.''}, National Museum of Mathematics, Pantograph ingredient/directions/use paragraphs (first appeared October 8, 2012): two midpoint long bars, two half-length bars, fixed endpoint, double/half tracing. See the \href{https://momath.org/mathmonday/math-monday-01-03-11/math-monday-linkages-the-world-b-x/}{MoMath article} (full locator in SOURCES.md). The millimeter kit, particular inserts, operating regions, anchor comparison and proof presentation are locally authored adaptations. No source photograph establishes this kit's physical readiness.

Givental, Nemirovskaya and Zakharevich, \emph{Math Circle by the Bay}, Preface viii--x, supports deep revisitable themes, manipulatives and varying pace. The particular routes, timing and hardware estimates here are editorial judgments, not results reported by that book. The recovered Week85 plan and production record preserve provenance and review history.

\sub{After a visit}
Record which packet/problem/inserts children actually used, whether they controlled a point independently, what comparisons they noticed, where the mechanism stuck, and which question they wanted next. Separate those observations from hypotheses about motor control or explanations. Save one identified overlay and a picture of the working assembly for a return visit. Prepared pages alone are not taught material.

All outputs remain unpiloted; exact fit, operation and hand clearance are physically untested. Digital review and rebuilding check the published artifacts, not the real equipment. The canonical Git publication does not update a Drive or classroom copy automatically.

\end{document}
